% sec-00-cover-letter.tex - IND Cover Letter (full)
\section*{Cover Letter}
\addcontentsline{toc}{section}{Cover Letter}

\noindent ChemicalQDevice \\
San Diego, California \\
July 1, 2026 \\[0.5em]
Food and Drug Administration \\
Center for Drug Evaluation and Research \\
Division of Oncology 1 \\
5901 to 5913 Ammendale Road \\
Beltsville, Maryland 20705 to 1266 \\[0.5em]

\noindent\textbf{Re: Initial Investigational New Drug Application, Serial 0000 - Daraxonrasib (RMC-6236), Phase 1, LLM-Directed Robotic Whipple in KRAS-Mutated PDAC}

\noindent Dear Division of Oncology 1 Reviewer:

On behalf of ChemicalQDevice, and in my capacity as sponsor-investigator, I am submitting this initial Investigational New Drug Application (IND v1.0) under 21 CFR \S 312.23 for a Phase 1, first-in-human, prospective open-label single-arm study of daraxonrasib (RMC-6236), an oral once-daily RAS(ON) multi-selective (pan-RAS) small-molecule inhibitor, administered in the perioperative setting of a large language model (LLM) directed eight-arm robotic pancreaticoduodenectomy (Whipple) for patients with resectable or borderline-resectable, KRAS G12 mutated pancreatic ductal adenocarcinoma (PDAC) \cite{daraxwhipple,onpremwhippl}. This application is assigned serial number 0000 as the initial submission and is transmitted with FDA Form 1571 as the cover sheet for the submission and the sponsor agreement, together with FDA Form 3674 (clinicaltrials.gov certification) and FDA Form 1572 \cite{cfr312paiuni}.

This submission is combined in nature. It is concurrently an IND under 21 CFR Part 312 for the investigational drug daraxonrasib and an Investigational Device Exemption (IDE) under 21 CFR Part 812 for the significant-risk, Class II collaborative on-premises robotic surgical device that directs the operative procedure as an early-feasibility study \cite{onpremwhippl,paiautospons}. The LLM functions strictly as a hash-pinned, second-opinion advisory oracle isolated outside the vendor kinematic stack; full device autonomy is prohibited in accordance with 21 CFR \S 312.21(e), and a verification-before-generation ten-gate suite governs every advisory output consistent with H.R. 9510, the Verification Before Generation in Physical AI Oncology Trials Act of 2026 \cite{congress9510,clintrustpai}. The dual regulatory pathways are presented in an integrated dossier so that the drug and device safety considerations are reviewed together.

We respectfully request the standard 30-day safety review period under 21 CFR \S 312.40, after which, absent a clinical hold, the sponsor-investigator intends to initiate the study at a single academic site \cite{phase1guidance}. The dose-escalation design follows a 3+3 schema across three oral once-daily dose levels (DL1 160 mg, DL2 220 mg, DL3 300 mg, the RASolute 302 recommended Phase 2 dose) with a 28-day dose-limiting toxicity window, enrolling up to 18 treated subjects \cite{daraxwhipple}.

\noindent\textbf{Contents enclosed.} The following components are submitted in the ReGARDD IND order, with the Cover Letter and FDA Forms 1571 and 3674 preceding the generated Table of Contents:

\noindent\begin{tabularx}{\textwidth}{L{2.2cm} Y}
\toprule
\textbf{Item} & \textbf{Component} \\
\midrule
1 & FDA Forms 1571 and 3674 (3674 Box B selected; PHS Act 402(j) registration requirements addressed) \\
2 & Table of Contents (generated) \\
3 & Introduction, including drug name, pharmacological class, structure, formulation, route, and duration \\
4 & General Investigational Plan, including the 3+3 escalation rationale and first-year plan \\
5 & Investigator Brochure \\
6 & Proposed Clinical Research (study protocol v1.0.0, informed consent, Form 1572, financial disclosures) \\
7 & Chemistry, Manufacturing, and Control information (drug substance, drug product, placebo, labeling, environmental assessment) \\
8 & Pharmacology and Toxicology information, including the integrated computational evidence base \\
9 & Previous Human Experience (RASolute 302, RASolve 301, Breakthrough Therapy designation) \\
10 & Additional Information (drug dependence, radioactive drugs, pediatric studies) \\
11 & Relevant Information and References with Back Matter \\
\bottomrule
\end{tabularx}
\figcaption{Table 1. Enclosed components of the initial IND, in ReGARDD order, with the Cover Letter and FDA Forms preceding the generated Table of Contents.}

\noindent\textbf{Sponsor-investigator and contact.} The sponsor-investigator is Kevin Kawchak, Chief Executive Officer, ChemicalQDevice, San Diego, California, who may be reached at kevink@chemicalqdevice.com. The sponsor-investigator's financial interest is disclosed under 21 CFR Part 54 and mitigated by a three-tier independent oversight structure comprising a Data Safety Monitoring Board with halt authority, an Independent Safety Monitor with real-time per-procedure stop authority, and a Physical AI Safety Review Committee operating on a cadence of no more than 90 days \cite{paisitedocpk}.

\noindent\textbf{AI-assisted assembly and real-time monitorability.} This IND was assembled under a single master prompt that executed a four-stage mermaid to draft to full to final build, with each generated file committed and pushed to one continuously updated GitHub pull request in real time, so that the complete authoring history of the submission is human-monitorable as it is produced and is reproducible from a hash-pinned deterministic commit (repository v4.3.0) \cite{oncotrialllm,natlpaiplatf}. Figure 1 shows the single-prompt build pipeline and Figure 2 shows the ReGARDD Table-of-Contents architecture that maps each numbered IND section to exactly one source file. The repository is available at \url{https://github.com/kevinkawchak/physical-ai-oncology-trials/tree/main/trial-ind}.

\begin{mermaidfig}
\node[mmgoal] (mp) at (0,0) {Master prompt\\prompts/prompt-ind.md};
\node[mmproc] (pa) at (4.6,0) {Process A\\write sub-prompts 1 to 4};
\node[mmin] (s1) at (0,-2.4) {Stage 1 mermaid\\22 grayscale figures};
\node[mmin] (s2) at (4.6,-2.4) {Stage 2 draft-ind\\12 sec tex scaffold\\+ instructions};
\node[mmaccent] (s3) at (9.2,-2.4) {Stage 3 full-ind\\full prose + 22 TikZ\\+ full-width tables};
\node[mmgoal] (s4) at (9.2,-5.0) {Stage 4 final-ind\\senior-author polish\\+ Overleaf zip};
\node[mmdark] (rel) at (4.6,-5.0) {Release v4.3.0\\README + CHANGELOG\\+ output-ind};
\node[mmctx] (gh) at (0,-5.0) {GitHub branch\\real-time auto-commit\\and auto-PR};
\draw[mmedge] (mp) -- (pa);
\draw[mmedge] (pa) to[out=-90,in=90,looseness=0.6] (s1.north);
\draw[mmedge] (s1) -- (s2);
\draw[mmedge] (s2) -- (s3);
\draw[mmedge] (s3) -- (s4);
\draw[mmedge] (s4) -- (rel);
\draw[mmedged] (s1) -- (gh);
\draw[mmedged] (s2.south) to[out=-90,in=0,looseness=0.6] (gh.east);
\draw[mmedged] (s3.south) to[out=-90,in=20,looseness=0.7] (gh.east);
\draw[mmedged] (s4.south) to[out=-90,in=-20,looseness=0.7] (gh.south);
\draw[mmedged] (gh.north) to[out=90,in=180,looseness=0.7] (s4.west);
\end{mermaidfig}
\figcaption{Figure 1. The single-prompt mermaid to draft to full to final build pipeline for the Phase 1 PDAC IND. One master prompt drives Process A (write the four sub-prompts) and then Process B (execute them as Stages 1 to 4). Dashed edges show that every generated file is committed and pushed to one continuously updated pull request in real time, and the feedback edge shows author review returning into the final stage, all without manual intervention.}

\begin{mermaidfig}
\node[mmaccent] (cl) at (0,0) {Cover Letter\\sec-00};
\node[mmproc] (f) at (0,-1.8) {1. FDA Forms 1571 / 3674\\sec-01};
\node[mmctx] (toc) at (0,-3.6) {2. Table of Contents\\generated};
\node[mmin] (i) at (0,-5.4) {3. Introduction\\sec-02 (3.1 to 3.5)};
\node[mmin] (gip) at (4.6,0) {4. General Investigational Plan\\sec-03 (4.1 to 4.7)};
\node[mmin] (ib) at (4.6,-1.8) {5. Investigator Brochure\\sec-04};
\node[mmin] (pcr) at (4.6,-3.6) {6. Proposed Clinical Research\\sec-05 (6.1 to 6.3)};
\node[mmin] (cmc) at (4.6,-5.4) {7. Chemistry, Manufacturing, Control\\sec-06 (7.1 to 7.2)};
\node[mmin] (pt) at (9.2,0) {8. Pharmacology / Toxicology\\sec-07 (8.1)};
\node[mmin] (phe) at (9.2,-1.8) {9. Previous Human Experience\\sec-08 (9.1 to 9.4)};
\node[mmin] (ai) at (9.2,-3.6) {10. Additional Information\\sec-09 (10.1 to 10.5)};
\node[mmin] (ri) at (9.2,-5.4) {11. Relevant Information\\sec-10};
\node[mmgoal] (bm) at (13.4,-2.7) {References + Back Matter\\sec-11};
\node[mmgoal] (main) at (13.4,-5.0) {main.tex\\assembles all};
\draw[mmedge] (cl) -- (f);
\draw[mmedge] (f) -- (toc);
\draw[mmedge] (toc) -- (i);
\draw[mmedge] (i.east) to[out=0,in=180,looseness=0.6] (gip.west);
\draw[mmedge] (gip) -- (ib);
\draw[mmedge] (ib) -- (pcr);
\draw[mmedge] (pcr) -- (cmc);
\draw[mmedge] (cmc.east) to[out=0,in=180,looseness=0.6] (pt.west);
\draw[mmedge] (pt) -- (phe);
\draw[mmedge] (phe) -- (ai);
\draw[mmedge] (ai) -- (ri);
\draw[mmedge] (ri.east) to[out=0,in=180,looseness=0.6] (bm.west);
\draw[mmedged] (main) to[out=160,in=-20,looseness=0.7] (cl.east);
\draw[mmedged] (main) -- (bm);
\draw[mmedged] (main) to[out=180,in=0,looseness=0.6] (toc.east);
\end{mermaidfig}
\figcaption{Figure 2. The ReGARDD IND Table of Contents reproduced as the file architecture of the build. The Cover Letter and the FDA Forms 1571 and 3674 precede the generated Table of Contents; the eleven numbered IND sections and the References and Back Matter each become exactly one sections/sec tex file assembled by main.tex, so the document structure and the repository structure match one-to-one.}

\noindent We appreciate the Division's review and welcome any questions during the 30-day period. The sponsor-investigator is available for a teleconference at the Division's convenience and will respond promptly to any request for additional information.

\noindent Respectfully submitted, \\[1.0em]
Kevin Kawchak \\
Chief Executive Officer and Sponsor-Investigator \\
ChemicalQDevice, San Diego, California \\
kevink@chemicalqdevice.com \\
IND v1.0; repository v4.3.0
